% chapters/02/section-04.tex
\section{Simulasi Fisika Diskret dan Integrasi Numerik}
\label{sec:simulasi-diskret}

Gerak pada permainan tampak berlangsung terus-menerus, padahal komputer
menghitungnya secara terputus. Bagian ini menjelaskan bagaimana gerak yang
seolah bersinambung itu dihasilkan dari perhitungan yang hanya dilakukan
sesekali, beserta harga yang harus dibayar atas cara tersebut.

\subsection{Simulasi Kontinu pada Mesin Diskret}

Sebuah bola yang jatuh bergerak menurut hukum gerak yang berlaku pada setiap
saat. Posisinya pada saat setengah detik, seperempat detik, atau sepersekian
detik mana pun dapat dihitung. Kenyataan itu disebut simulasi kontinu.

Komputer tidak dapat menghitung pada setiap saat karena jumlah saat tidak
terhingga. Yang dapat dilakukannya adalah menghitung posisi pada beberapa saat
tertentu, misalnya setiap seperenam puluh detik, lalu menghubungkan hasilnya
menjadi gerak yang tampak mulus. Selisih antarwaktu perhitungan itulah yang
disebut panjang langkah, dilambangkan $\Delta t$, dan sama dengan selang waktu
yang telah Anda kenal pada Bagian~\ref{sec:model-waktu}.

Karena itu, hasil simulasi pada komputer selalu merupakan hampiran atas gerak
yang sebenarnya. Pertanyaan yang tepat bukanlah ``apakah hasilnya tepat?'',
melainkan ``seberapa besar selisihnya, dan apakah selisih itu masih dapat
diterima untuk keperluan permainan?'' \cite{catto-integration}.

\subsection{Kaidah Euler}

Cara paling sederhana untuk memajukan simulasi adalah kaidah Euler, yang
menghitung nilai baru berdasarkan nilai lama dan laju perubahannya. Untuk
keadaan gerak, kaidah ini diterapkan pada dua tingkat secara berurutan.

\begin{align}
v &\leftarrow v + a\,\Delta t, \label{eq:euler-v}\\
p &\leftarrow p + v\,\Delta t. \label{eq:euler-p}
\end{align}

Persamaan~\eqref{eq:euler-v} memperbarui kecepatan menggunakan percepatan,
kemudian Persamaan~\eqref{eq:euler-p} memperbarui posisi menggunakan kecepatan
yang baru saja diperbarui. Urutan tersebut tidak boleh ditukar. Apabila posisi
diperbarui lebih dahulu dengan kecepatan yang lama, hasilnya sama dengan
memakai kecepatan satu langkah sebelumnya. Pada benda yang sedang jatuh,
kesalahan seperti itu menumpuk ke satu arah, sehingga benda terlihat jatuh
lebih cepat daripada seharusnya.

Bentuk yang dipakai pada buku ini, yaitu memperbarui kecepatan lebih dahulu,
dikenal sebagai kaidah Euler semi-implisit. Kelebihannya bukan pada
ketelitiannya, melainkan pada sifatnya yang lebih stabil dan mudah dipahami,
sehingga banyak dipakai pada permainan.

\subsection{Galat dan Panjang Langkah}

Seberapa besar selisih yang timbul? Gambar~\ref{fig:integrasi-euler} memuat
jawabannya untuk satu contoh sederhana: sebuah benda dilepaskan dari
ketinggian 10 satuan tanpa kecepatan awal, dengan percepatan gravitasi 10
satuan per detik kuadrat.

\begin{figure}[htbp]
\centering
% chapters/02/figures/integrasi-euler.tikz
% Benda jatuh bebas: lintasan eksak dibandingkan dua langkah Euler.
% y0 = 10 satuan, v0 = 0, g = 10 satuan/detik^2.
\begin{tikzpicture}
\begin{axis}[
  width=0.92\textwidth,
  height=6.4cm,
  xlabel={waktu $t$ (detik)},
  ylabel={ketinggian $y$ (satuan)},
  xmin=0, xmax=2.05, ymin=-17, ymax=12,
  axis lines=left,
  grid=both,
  grid style={ubLine!45},
  tick label style={font=\scriptsize},
  label style={font=\scriptsize},
  legend pos=north east,
  legend style={font=\scriptsize, fill=white, draw=ubLine}
]
\addplot[domain=0:2, samples=90, ubPrimary, very thick]
  {10 - 5*x^2};
\addlegendentry{lintasan eksak $y = 10 - 5t^2$}

\addplot[ubRed, thick, dashed, mark=*, mark size=1.5pt]
  coordinates {(0,10)(0.5,7.5)(1,2.5)(1.5,-5)(2,-15)};
\addlegendentry{Euler, $\Delta t = 0{,}5$}

\addplot[ubGreen, thick, dashdotted, mark=square*, mark size=1.2pt]
  coordinates {(0,10)(0.25,9.375)(0.5,8.125)(0.75,6.25)(1,3.75)
               (1.25,0.625)(1.5,-3.125)(1.75,-7.5)(2,-12.5)};
\addlegendentry{Euler, $\Delta t = 0{,}25$}
\end{axis}
\end{tikzpicture}

\caption{Lintasan benda jatuh bebas: perbandingan hasil eksak dan dua panjang langkah Euler}
\label{fig:integrasi-euler}
\end{figure}

Setelah dua detik, lintasan yang sebenarnya menempatkan benda pada ketinggian
$-10$ satuan. Dengan panjang langkah $0{,}5$ detik, kaidah Euler menempatkannya
pada $-15$ satuan; dengan panjang langkah $0{,}25$ detik, pada $-12{,}5$
satuan. Selisihnya menyusut dari 5 satuan menjadi 2,5 satuan ketika panjang
langkah diperkecil menjadi separuh.

Pola tersebut berlaku umum pada kaidah Euler: \textbf{galat berbanding lurus
dengan panjang langkah}. Memperkecil panjang langkah menjadi separuh
memperkecil galat menjadi kira-kira separuh pula. Sifat ini disebut ketelitian
tingkat pertama. Konsekuensinya praktis, yaitu menyetel panjang langkah adalah
cara paling langsung untuk menukar beban perhitungan dengan ketelitian.

\begin{catatabox}
Galat pada contoh tersebut bukan kesalahan penulisan
program, melainkan akibat tak terhindarkan dari menghampiri gerak yang
bersinambung dengan perhitungan yang terputus. Menambah panjang langkah pada
umumnya lebih merugikan daripada menguntungkan, dan menghaluskan tampilan
tidak memperbaiki perhitungan yang mendasarinya.
\end{catatabox}

\subsection{Kestabilan dan Pilihan Praktis}

Selain kurang teliti, langkah yang terlalu panjang dapat membuat simulasi
tidak stabil. Objek yang bergerak cepat dapat melewati seluruh lebar dinding
dalam satu langkah, sehingga pemeriksaan tabrakan yang dilakukan setelah
langkah tersebut tidak menemukan apa pun dan objek keluar dari ruang
permainan. Gejala ini disebut penembusan (\istilah{tunneling}) dan akan Anda
jumpai kembali pada Bab~VII.

Tiga cara penanganan berikut lazim dipakai, dan ketiganya dapat digabungkan.

\begin{enumerate}
  \item \textbf{Membatasi panjang langkah maksimum.} Apabila selang waktu satu
        bingkai ternyata jauh lebih lama daripada biasanya, nilainya dipangkas
        sebelum dipakai. Cara ini mencegah simulasi melompat jauh ketika
        komputer sesaat tersendat.
  \item \textbf{Memecah langkah menjadi beberapa sub-langkah.} Satu bingkai
        yang panjang dibagi menjadi beberapa langkah pendek yang dijalankan
        berurutan, sehingga perpindahan pada setiap langkah tetap kecil.
  \item \textbf{Memakai langkah tetap.} Simulasi selalu maju dengan panjang
        langkah yang sama, sebagaimana telah diuraikan pada
        Bagian~\ref{sec:model-waktu}.
\end{enumerate}

Pada Blok~I dan Blok~II, ketiga cara itu Anda terapkan sendiri pada kode Anda.
Pada Blok~III, Unity menyediakan langkah tetap melalui tahap pembaruannya
sendiri, dan tugas Anda berpindah menjadi memilih nilai panjang langkah yang
tepat serta mengenali gejala ketika nilai tersebut kurang sesuai. Perpindahan
tanggung jawab semacam inilah yang dibahas lebih jauh pada
Bagian~\ref{sec:pustaka-mesin}.
